\documentclass[11pt]{article}

\usepackage[utf8]{inputenc}
\usepackage[T1]{fontenc}
\usepackage{amsmath, amssymb}
\usepackage{geometry}
\usepackage{booktabs}
\usepackage{listings}
\usepackage{xcolor}
\usepackage{hyperref}
\usepackage{enumitem}
\usepackage{tikz}
\usepackage{bm}
\usetikzlibrary{positioning, arrows.meta}

\geometry{margin=1in}

\definecolor{codebg}{rgb}{0.96,0.96,0.96}
\lstset{
  backgroundcolor=\color{codebg},
  basicstyle=\ttfamily\small,
  breaklines=true,
  frame=single,
  language=Python,
  showstringspaces=false,
  columns=fullflexible,
  keywordstyle=\color{blue!70!black}
}

\newcommand{\diver}{\nabla\!\cdot}
\newcommand{\grad}{\nabla}
\newcommand{\bx}{\mathbf{x}}
\newcommand{\bp}{\mathbf{p}}
\newcommand{\R}{\mathbb{R}}

\title{RBF-FD Solver: Least Squares Method for Operator Reproduction}
\author{Dennis Corraliza}
\date{9/2/2026}

\begin{document}
\maketitle

\pagebreak

\section{Problem Setting:}

Following the standard RBF-FD methods, let $u \in \mathcal{F}$ be a solution for a PDE in form of:
\begin{equation}
	- Lu = -\diver\!\big(A \grad u\big) = f \quad \text{in } \Omega,
	\label{eq:PDE}
\end{equation}

Let $1<p<\infty$, $A \in C^{0,1} (\bar{\Omega}; \mathbb{R}^{d \times d}):\bar{\Omega} \rightarrow \mathbb{R}^{d \times d}$ is an symmetric positive definite anisotropic tensor, where
\begin{equation}
	\eta |\xi|^2 \le \bf{A}(\bf{x}) \xi \cdot \xi \le H |\xi|^2 \text{, } 0 < \eta \le H
\end{equation}

and $f \in L^p(\Omega)$ is the forcing term acting in the domain $\Omega$. The problem is subject to Neumann and Dirichlet Boundary Conditions partitioned within $\partial \Omega$ of the form:
\begin{align}
  u &= g_D \quad \text{on } \Gamma_D \subset \partial\Omega, \\
  \frac{\partial u}{\partial n} \equiv n \cdot A \grad u &= g_N \quad
  \text{on } \Gamma_N \subset \partial\Omega,
  \label{eq:neumann-bc}
\end{align}

For simplicity, we assume $\partial \Omega = \Gamma_D \cup \Gamma_N$ and $\Gamma_D \cap \Gamma_N = \emptyset$. Furthermore, let $g_D, g_N \in W^{k-1,p}(\partial \Omega)$. By \textit{Agmon-Douglis-Nirenberg} Theorem, we have the following estimate for $u$:
\begin{equation}
\|u\|_{W^{k,p}(\Omega)} \leq C\left(\|Lu\|_{L^p(\Omega)} \sum_{j=1}^{m} \|g_i\|_{W^{k-1,p}(\partial\Omega)} + \|u\|_{L^p(\Omega)} \right).
\end{equation}

\section{RBF-FD Method:}

\subsection{Local Approximation setup}
To obtain an approximation of the operator $L$ at point $\bf{z}$, we can approximate $\hat{u} \approx u$ on a neighborhood $\Omega \subseteq \mathbb{R}^n$ of $\bf{z}$ such that we approximate $Lu \approx L\hat{u}$. We define the radial basis functions in the form $\phi_i(\bm{x}) := \phi(||\bm{x}-\bm{y}_i||): \Omega \rightarrow \mathbb{R}$ for $i = 1, \hdots , m$ and centers $\left\{ \bm{y}_i \right\}_{i=1}^m \subseteq \Omega$. Define the vector space $V_Y = span\left\{ \phi_1, \dots, \phi_m \right\}$. Let $X = \left\{ \bm{x}_1, \hdots, \bm{x}_n \right\} \subseteq \Omega$ be a set of stencil nodes for a point $\bf{z}$ in $\Omega$. For the classical approximation scheme, let $X=Y$ and $m=n$. We can now define $\bm{\Phi} \in \mathbb{R}^{n \times m}$, where $\bm{\Phi}_{ij} = \phi_j(\bm{x}_i)$. The role of these basis functions will be to approximate the operator $L$. The approximation of $u$ takes the form of:
\begin{equation}
	\hat{u}(\bf{x}) =  \sum_{i=1}^n \lambda_i \phi_i(\bf{x})
\end{equation}

This will lead to the form $\bf{\Phi} \bf{\lambda} = \hat{\bf{u}}$ for $\hat{\bf{u}} = [ u(\bf{x_1}), \hdots, u(\bf{x_n}) ]^T.$ We can simply for $\lambda$ to obtain:
\begin{equation}
\hat{u}(\bf{x}) = \sum_{i=1}^n (\bf{\Phi}^{-1} \bf{u})_i \phi_i(\bf{x}) = \bf{\phi}(\bf{x})^T \bf{\Phi}^{-1} \bf{u}
\end{equation}

The approximation of $Lu$ at $z$ is in the form of:
\begin{equation}
	Lu (\bf{z}) \approx L \hat{u} (\bf{z}) = L \bf{\phi}^T(\bf{z}) \bf{\lambda} = L \bf{\phi}^T(\bf{z}) \bf{\Phi}^{-1} \bf{u} := w^{L, X}(\bf{z}) \bf{u}
\end{equation}

Thus, approximation the operator turns into local solver for the weights $w_i$ from the form:
\begin{equation}
	\bf{\Phi}^T \bf{w^{L, X}(\bf{z})} = L \bf{ \phi} (\bf{z})
\end{equation}

Through the method of undefined coefficients on the basis $\{\phi_i\}$ we get the same resulting set of weights. Finally, by defining the basis for $V_Y$ as radial basis functions results in the Radial Basis Function generated Finite Difference method (RBF-FD). 

\subsection{Polynomial Augmentation}

For conditionally positive definite kernels, we augment them through the form:

\section{Weight Constraint RBF-FD Setup:}

Rather than assuming a form for the solution, our goal is to approximate the operator through weights dual to $u$ sampled at nodes. Suppose $m > n$, thus more center nodes than stencil nodes per stencil. Additionally, instead of assuming exactness of the form $L\phi_j(\bm{x}) = \sum_i^n w_i(\bm{x}) \phi_j(\bm{x}_i)$ for $j = 1, \hdots, m$ as in the classical RBF-FD, we define a residual function:
\begin{align}
	r_j(\bm{x}) :&= L\phi_j(\bm{x}) - \sum_{i=1}^n w_i(\bm{x}) \phi_j(\bm{x_i})\\
	r(\bm{x}) &= L\bm{\phi}(\bm{x}) - \bm{\Phi^T} \bm{w}(\bm{x}), \text{ } \bm{w}: \mathbb{R}^n \rightarrow \mathbb{R}^n \text{ for } \bm{x} \in \Omega.
\end{align}

Suppose $\left\{p_i\right\}_{i=1}^l$ is a set of polynomials augmenting the radial basis functions. Let $\Pi_k=span\left\{p_i\right\}_{i=1}^l$ with polynomial degree up to $k$. We want to reproduce the operator exactly on the polynomial basis: 
\begin{align}
	Lp_j(\bm{x}) :&= \sum_{i=1}^n w_i(\bm{x}) p_j(\bm{x_i}) \text{ for } j = 1, \hdots, l\\
	L\bm{p}(\bm{x}) &= \bm{P^T} \bm{w}
\end{align}

where $\bm{P}_{ij} := p_j(\bm{x}_i) \in \mathbb{R}^{n \times l}$. Thus, we have a constrained least squares problem in the form:
\begin{equation}
	\left\{
	\begin{aligned}
		\min J(\bm{w}(\bm{x})) &:= \min \frac{1}{2}\|\bm{r}(\bm{x})\|_2^2 \quad \text{for } \bm{x} \in \Omega \\
		\text{subject to } \bm{\Phi}^T \bm{w} &= L\bm{p}
	\end{aligned}
	\right.
\end{equation}

We can solve this setup by defining the Lagrangian $\mathcal{L}$ for the problem as follows:
\begin{align*}
	\mathcal{L}(\bm{w}, \bm{\nu}) :&= J(\bm{w}) + <\bm{\nu}, \bm{P}^T \bm{w} - L\bm{p}(\bm{x}) > \\
	&= \frac{1}{2}<\bm{\Phi}^T \bm{w}, \bm{\Phi}^T \bm{w}> - <\bm{\Phi}^T \bm{w}, L\bm{\phi}> + \frac{1}{2}<L\bm{\phi}, L\bm{\phi}>  + <\bm{\nu}, \bm{P}^T \bm{w} - L\bm{p} >
\end{align*}

\section{Consistency Conditions for Finite Spaces and Interpolation Results}

The resulting KKT equations for this system are then:
\begin{align}
	\nabla_w \mathcal{L} &= 0: \bm{\Phi} L\bm{\phi} = \bm{\Phi} \bm{\Phi}^T \bm{w} + \bm{P} \bm{\nu} \label{eq:normal_eq}\\
	\nabla_\nu \mathcal{L} &= 0: \bm{P}^T \bm{w} = L\bm{p}
	\label{eq:constraint}
\end{align}

The augmented matrix then takes the form:
\begin{equation}
	\begin{bmatrix}
		\bm{\Phi} \bm{\Phi}^T & \bm{P} \\
		\bm{P}^T & \bm{0}
	\end{bmatrix}
	\begin{bmatrix}
		\bm{w} \\ \bm{\nu}
	\end{bmatrix} =
	\begin{bmatrix}
		\bm{\Phi} L\bm{\phi} \\ L\bm{p}
	\end{bmatrix}
\end{equation}

The primary goal is to approximate $Lu(\bm{x}) \approx \bm{w}^T(\bm{x}) \bm{u}$, with $\bm{u}=\left[ u(\bm{x}_1), \hdots, u(\bm{x}_n) \right]^T$. Define $L_h f(\bm{x}) := \bm{w}^T (\bm{x}) \bm{f}$ for some function $f$. Thus, we have the following:
\begin{align*}
	L_h f(\bm{x}) &= \bm{w}^T(\bm{x}) \bm{f}\\
			&= \begin{bmatrix}
					\bm{w}^T & \bm{\nu}^T
				\end{bmatrix}
				\begin{bmatrix}
					\bm{f} \\ \bm{0}
				\end{bmatrix}\\
			&= \left(
	\begin{bmatrix}
		\bm{\Phi} \bm{\Phi}^T & \bm{P} \\
		\bm{P}^T & \bm{0}
	\end{bmatrix}^{-1}
	\begin{bmatrix}
		\bm{\Phi} L\bm{\phi} \\ L\bm{p}
	\end{bmatrix}
				\right)^T 
				\begin{bmatrix}
					\bm{f} \\ \bm{0}
				\end{bmatrix}\\
			&= \begin{bmatrix}
					L\bm{\phi}^T \bm{\Phi}^T & L\bm{p}^T
				\end{bmatrix}
	\begin{bmatrix}
		\bm{\Phi} \bm{\Phi}^T & \bm{P} \\
		\bm{P}^T & \bm{0}
	\end{bmatrix}^{-T}
				\begin{bmatrix}
					\bm{f} \\ \bm{0}
				\end{bmatrix}\\
			&= \begin{bmatrix}
					L\bm{\phi}^T \bm{\Phi}^T & L\bm{p}^T
				\end{bmatrix}
				\begin{bmatrix}
					\bm{q} \\ \bm{\alpha}
				\end{bmatrix}\\
			&= L\bm{\phi}^T (\bm{x}) \left(\bm{\Phi}^T \bm{q}\right) + L\bm{p}^T (\bm{x}) \bm{\alpha}
\end{align*}

Next, we can show that the coefficients $\bm{q}, \bm{\alpha}$ relate to the coefficients of $u$ as follows:

\begin{equation}
	\begin{bmatrix}
		\bm{f} \\ \bm{0}
	\end{bmatrix} = 
	\begin{bmatrix}
		\bm{\Phi} \bm{\Phi}^T & \bm{P} \\
		\bm{P}^T & \bm{0}
	\end{bmatrix}
	\begin{bmatrix}
		\bm{q} \\ \bm{\alpha}
	\end{bmatrix}
\end{equation}

Define $V_h = \left\{ \rho(x) : \rho(x) = \sum_{i=1}^m \lambda_i \phi_i(x) \text{ such that } \lambda \in Range\left(\bm{\Phi}^T \right) \right\} \subseteq V_Y$. Let $s \in V_h + \Pi_k$. Thus, $s(x) = \bm{\phi}^T(\bm{x})\bm{\Phi}^T q + \bm{p}^T(\bm{x}) \bm{\alpha}$ where $\bm{q} \in \mathbb{R}^{n}$ and $\bm{\alpha} \in \mathbb{R}^{l}$. Then:
\begin{align}
	\bm{s} = s(X) &= \bm{\Phi} \bm{\Phi}^T \bm{q} + \bm{P} \bm{\alpha} = f(X) = \bm{f}\\
	\bm{P}^T \bm{q} &= \bm{0}
\end{align}

Hence, we get interpolation of $f$ at the node set $X$. Furthermore, from the KKT matrix we can show that $\bm{\Phi} \bm{r} = \bm{P} \bm{\nu}$ which implies that $\bm{\Phi} \bm{r} \in range(\bm{P}) = \left( ker \left(\bm{P}^T \right) \right)^\perp$. Thus, we get the relationship:
\begin{equation}
\bm{r}^T \left( \bm{\Phi}^T \bm{z} \right)=\bm{0} \text{ for } \bm{P}^T \bm{z} = \bm{0} \label{eq:res_orth}
\end{equation}

Finally, we check the consistency error for $s \in V_h+\Pi_k$ with $\bm{P}^T \bm{q}=\bm{0}$.
\begin{align*}
	Ls(x)-L_hs(x) &= L\left(\bm{\phi}^T(\bm{x})\bm{\Phi}^T \bm{q} + \bm{p}^T(\bm{x}\right) \bm{\alpha}) - \bm{w}^T(\bm{x})\bm{s}\\
	&= L\bm{\phi}^T(\bm{x})\bm{\Phi}^T \bm{q} + L\bm{p}^T(\bm{x}) \bm{\alpha} - \bm{w}^T(\bm{x})(\bm{\Phi} \bm{\Phi}^T \bm{q} + \bm{P} \bm{\alpha})\\
	&= L\bm{\phi}^T(\bm{x})\bm{\Phi}^T \bm{q} + L\bm{p}^T(\bm{x}) \bm{\alpha} - \bm{w}^T(\bm{x})\bm{\Phi} \bm{\Phi}^T \bm{q} - \bm{w}^T(\bm{x}) \bm{P} \bm{\alpha}\\
	&= \left( L\bm{\phi}(\bm{x}) - \bm{\Phi}^T \bm{w}(\bm{x}) \right)^T \bm{\Phi}^T \bm{q} + \left( L\bm{p}^T(\bm{x}) - \bm{P}^T \bm{w}(\bm{x}) \right)^T \bm{\alpha}\\
	&= \bm{r}^T \bm{\Phi}^T \bm{q} + \left( L\bm{p}(\bm{x}) - \bm{P}^T \bm{w}(\bm{x}) \right)^T \bm{\alpha}\\
\end{align*}
\begin{equation}
	Ls(x)-L_hs(x) = 0 \text{ due to (\ref{eq:constraint}) and (\ref{eq:res_orth}}) \label{eq:repro}
\end{equation}

\section{Consistency Results for Function Approximations}.

Let $u$ be a solution to (\ref{eq:PDE}) and $s \in V_h+\Pi_k$. We can say $u = s + (u - s)$ to get the following:
\begin{align*}
\mathcal{E}[u](\bm{x}) :&= Lu(\bm{x}) - L_hu(\bm{x})\\
	&= Lu(\bm{x}) - \bm{w}^T(\bm{x})\bm{u}\\
	&= L((u-s) + s)(\bm{x}) - \bm{w}^T(\bm{x})\bm{u} - L_h s(\bm{x}) + L_h s(\bm{x})\\
	&= L(u-s)(\bm{x}) + (L s(\bm{x}) - L_h s(\bm{x})) + L_h(s-u)(\bm{x})\\
	&= L(u-s)(\bm{x}) - \bm{w}^T(\bm{x})(\bm{u}-\bm{s}) \text{, by (\ref{eq:repro})}\\
	&= Le(\bm{x}) - \bm{w}^T(\bm{x})\bm{e} \text{, } e(\bm{x})=u(\bm{x})-s(\bm{x})\\
	&= \mathcal{E}[e](\bm{x})
\end{align*}

\begin{thebibliography}{2}

\bibitem{TrefethenBau}
Agmon, S., Douglis, A. and Nirenberg, L. (1959). \textit{Estimates near the boundary for solutions of elliptic partial differential equations satisfying general boundary conditions}. I. Comm. Pure Appl. Math., 12: 623-727. \\
\url{https://doi.org/10.1002/cpa.3160120405}

\bibitem{Slak}
Slak, J. (2026). \textit{Adaptive RBF-FD method}. PhD Thesis, University of Ljubljana.

\end{thebibliography}

\end{document}